% Generated draft: requires mathematical and source review.
\documentclass[11pt]{article}
\usepackage[margin=25mm]{geometry}
\usepackage{fontspec}
\setmainfont{Latin Modern Roman}
\usepackage{amsmath,amssymb,mathtools,unicode-math}
\setmathfont{Latin Modern Math}
\usepackage{enumitem,needspace,xurl,hyperref}
\setlength{\parindent}{0pt}
\setlength{\parskip}{5pt}
\setlength{\emergencystretch}{3em}
\newcommand{\N}{\mathbb{N}}
\newcommand{\Z}{\mathbb{Z}}
\newcommand{\Q}{\mathbb{Q}}
\newcommand{\R}{\mathbb{R}}
\newcommand{\C}{\mathbb{C}}
\newcommand{\F}{\mathbb{F}}
\newcommand{\A}{\mathbb{A}}
\newcommand{\PP}{\mathbb{P}}
\newcommand{\E}{\mathbb{E}}
\newcommand{\eps}{\varepsilon}
\newcommand{\dd}{\,\mathrm{d}}
\begin{document}
\section*{UnsolvedMath ID: 30003841}
\textbf{Source problem code:} catalogue entry
\subsection{Definitions and mathematical statement}
Let \(V_{\mathrm{br}}\) denote the braided Thompson group named in the supplied record. The record describes a \emph{closed spraige} as a closed split--braid--merge band diagram representing conjugacy data in \(V_{\mathrm{br}}\). This is a descriptive definition: the supplied material does not specify the admissible diagrams, their closure operation, or the equivalence relations and diagram moves governing them. It also does not supply a definition or presentation of \(V_{\mathrm{br}}\).

\textbf{Research question.} What further information about \(V_{\mathrm{br}}\) can be extracted from closed spraiges? In the supplied expanded formulation, the requested information is structural or dynamical information beyond their use in studying conjugacy.

\textbf{Unresolved scope.} The record does not designate a particular invariant, structural property, dynamical action, or recovery procedure. Nor does it specify whether recovery is to be effective, complete, or invariant under a particular equivalence of closed diagrams. Consequently, this is an exploratory research question, and the supplied evidence does not determine a single quantified conjecture. A precise specialization requires both a formal diagram model and a specified target of recovery. The supplied assessment mentions a conjugacy result but does not state it; no particular conjugacy classification theorem is asserted here.
\subsection{Short English statement}
The problem asks what closed spraiges reveal about the braided Thompson group \(V_{\mathrm{br}}\), beyond their role in studying conjugacy. The record suggests structural or dynamical information but specifies neither a particular invariant nor a formal recovery criterion. This statement-only draft therefore preserves the exploratory question and requires human mathematical review.
\subsection{Sources}
\begin{itemize}
\item Local UnsolvedMath registry, permanent ID \(30003841\), title \emph{Closed Spraiges in the Braided Thompson Group}, exact source problem code \texttt{catalogue entry}; collection: \emph{Oberwolfach Reports Open Problems}. The recorded proposed year is \(2018\); no individual proposer is supplied.
\item Supplied source citation: \emph{Cohomological and Metric Properties of Groups of Homeomorphisms of R} (2019), Oberwolfach Reports, pp.\ 1579--1633, DOI \texttt{10.4171/owr/2018/26}. \url{https://doi.org/10.4171/owr/2018/26}. The supplied background describes these reports as records of mathematical workshops and their open-problem discussions.
\item Supplied literature source, labeled \emph{Oberwolfach Report 2019}: \url{https://ems.press/content/serial-article-files/46748?nt=1}.
\item Literature assessment dated exactly \texttt{2026-08-21}: the supplied curation labels the problem open and describes it as a broad follow-up concerning further invariants of closed spraiges, which the cited conjugacy result does not settle. This is supplied dated curation, not an independently verified assessment; the record itself requests independent verification of time-sensitive status and cited claims.
\item The registry labels the statement \texttt{corrected\_verified} and reports that it was checked against the official Oberwolfach report, preserving its hypotheses, notation, and open target. These are supplied verification claims; no independent source inspection is asserted in this draft.
\end{itemize}
\end{document}
